% OWNER: theory-a
% STATUS: review (poprawiono: usunięto kopiuj-wklej disclaimer ComfyUI, uzupełniono cytowania Ulrich, Pasini, Pallasmaa, wygładzono styl)
% SOURCES: notes/SPIS_TRESCI_SZCZEGOLOWY.md §2; GLOSSARY.md; notes/HIPOTEZY_REWIZJA.md
\chapter{Projektowanie wnętrz, psychologia środowiskowa i doświadczenie przestrzeni}
\label{ch02:top}
\ChapterStatus{review}{theory-a}

Rozdział buduje teoretyczną i projektową ramę dla platformy IDA (\emph{AI Interior Design Dialogue Research Platform}). Postrzeganą regeneracyjność, biofilię oraz funkcję przestrzeni traktuję w nim jako fundamentalne kategorie projektowe, a nie wyłącznie jako abstrakcyjne konstrukty psychometryczne.
Celem rozdziału jest uzasadnienie decyzji metodycznych, które w toku pełnego doświadczenia (Full) prowadzą od etapu konfiguracji pomieszczenia (Room Setup) do syntezy wariantu łączącego nastrój z funkcją (\texttt{mixed\_functional}) w eksperymentalnej macierzy generacji \footcite{idaPlatform2026}.

\section{Środowiska regeneracyjne: ART i SRT wobec wnętrza mieszkalnego}
\label{ch02:sec:art}

Wizualizacja wnętrza --- również ta syntetyzowana przez modele sztucznej inteligencji --- pełni w praktyce projektowej rolę obrazu zapośredniczającego koncepcję. Nie zastępuje ona projektu budowlanego ani wiedzy wykonawczej.
W badaniach przez projektowanie (RtD) obraz ten staje się jednak nośnikiem hipotez o jakości doświadczenia przestrzennego: pozwala badać, czy proponowana atmosfera przybliża mieszkańca do pożądanego stanu ukojenia i regeneracji, czy stanowi jedynie powierzchowną dekorację kadru.

Teoria regeneracji uwagi (\emph{Attention Restoration Theory}, ART) sformułowana przez Stephena Kaplana opisuje warunki, w których środowisko fizyczne wspiera odzyskiwanie wyczerpanej uwagi ukierunkowanej. Kaplan wyróżnia cztery konstytutywne cechy środowiska regeneracyjnego: fascynację (\emph{fascination}), poczucie bycia poza codziennością (\emph{being away}), rozległość i spójność przestrzeni (\emph{extent}) oraz zgodność środowiska z intencjami człowieka (\emph{compatibility}) \footcite{kaplan1995restorative}.
Z kolei teoria redukcji stresu (\emph{Stress Recovery Theory}, SRT) kładzie nacisk na szybką, ewolucyjnie ukształtowaną odpowiedź afektywną na kontakt z elementami przyrody. Klasyczne badania Rogera Ulricha wykazały bezpośredni wpływ widoku z okna na proces rekonwalescencji pacjentów \footcite{ulrich1984window}, a późniejsze eksperymenty laboratoryjne potwierdziły psychofizjologiczne mechanizmy redukcji pobudzenia układu współczulnego pod wpływem bodźców naturalnych \footcite{ulrich1991stress}.

Dla projektanta wnętrz mieszkalnych obie te teorie niosą jednoznaczny wniosek: zdolność przestrzeni do regeneracji psychofizycznej nie jest przypadkowym dodatkiem, lecz mierzalną jakością relacji człowieka z otoczeniem. \textcolor{orange}{Jakość tę kształtuje się operowaniem światłem dziennym, taktylną obecnością naturalnych materiałów (drewno, kamień, len), zrównoważoną złożonością wizualną oraz czytelnością układu stref \footcite{pallasmaa2012oczy}.}
Platforma IDA nie zakłada naiwnie, że wygenerowany obraz sam w sobie „leczy” zmęczenie; wykorzystuje natomiast mapę nastroju do zidentyfikowania luki między stanem obecnym a pożądanym w danym pokoju.

\begin{figure}[htbp]
  \centering
  \Todo{wstaw plik: figures/ch02/ch02-restorative-framework.pdf --- schemat ART/SRT $\rightarrow$ kategorie projektowe wnętrza (światło, materiał, natura, złożoność, układ)}
  \caption[Od teorii środowisk regeneracyjnych do kategorii projektowych wnętrza. Opracowanie własne na podstawie Kaplana oraz Ulricha.]{Od teorii środowisk regeneracyjnych do kategorii projektowych wnętrza.
    Opracowanie własne na podstawie \textcite{kaplan1995restorative} oraz \textcite{ulrich1984window,ulrich1991stress}.}
  \label{fig:ch02:restorative-framework}
\end{figure}

\section{Postrzegana regeneracyjność (PRS) jako kategoria projektowa}
\label{ch02:sec:prs}

Terry Hartig i współpracownicy opracowali skalę postrzeganej regeneracyjności (\emph{Perceived Restorativeness Scale}, PRS), która przekłada wymiary teorii Kaplana na standaryzowany pomiar psychometryczny \footcite{hartig1997prs}. Skala ta doczekała się wielu adaptacji, w tym zoptymalizowanej wersji skróconej PRS-11 \footcite{pasini2014prs11}.

W interfejsie IDA pomiar ten nie przybiera jednak formy nużącej ankiety z pozycjami Likerta. Uczestnik zaznacza stan emocjonalny na dwuwymiarowej mapie nastroju (\emph{mood grid}), gdzie oś~$X$ rozpięta jest między biegunem \emph{energetyzującym} a \emph{uspokajającym}, zaś oś~$Y$ między nastrojem \emph{monotonnym} a \emph{inspirującym}.
Jest to świadoma adaptacja projektowa (tab.~\ref{tab:ch02:prs-mood-grid}): oś pozioma przybliża wymiar wyciszenia (\emph{being away}), oś pionowa odzwierciedla fascynację (\emph{fascination}), natomiast poczucie spójności (\emph{coherence}) ujawnia się w powiązaniu nastroju z funkcjonalną konfiguracją pokoju \footcite{idaPlatform2026}.

W architekturze IDA pojęcie PRS obejmuje trzy powiązane stany:
\begin{itemize}
  \item \textbf{ideal} (\texttt{prs\_ideal}): ogólny, pożądany klimat emocjonalny deklarowany w profilu użytkownika;
  \item \textbf{current} (\texttt{prs\_current}): odczuwany, zastany stan emocjonalny konkretnego pokoju;
  \item \textbf{target} (\texttt{prs\_target}): docelowy nastrój, jaki użytkownik pragnie osiągnąć w tej przestrzeni.
\end{itemize}

Różnica wektorowa między stanem \texttt{prs\_current} a \texttt{prs\_target} definiuje lukę nastrojową, stając się wytyczną dla generatora AI (pytanie~\RQ{5}). Badanie nie ocenia laboratoryjnego „przyrostu regeneracji” po ekspozycji na render, lecz sprawdza, czy użytkownicy preferują warianty uwzględniające ową metamorfozę nastroju.

\begin{table}[htbp]
  \centering
  \caption{Wymiary PRS/ART a operacjonalizacja mapy nastroju (mood grid) w IDA.}
  \label{tab:ch02:prs-mood-grid}
  \begin{tabular}{@{}>{\raggedright\arraybackslash}p{0.24\textwidth}>{\raggedright\arraybackslash}p{0.32\textwidth}>{\raggedright\arraybackslash}p{0.36\textwidth}@{}}
    \toprule
    Wymiar PRS / ART & Odpowiednik w mood grid IDA & Etap badania w platformie \\
    \midrule
    \emph{Being away} (oderwanie / spokój)
      & Oś~$X$: energetyzujący~$\leftrightarrow$~uspokajający
      & \texttt{prs\_ideal}, \texttt{prs\_current}, \texttt{prs\_target} \\
    \emph{Fascination} (fascynacja)
      & Oś~$Y$: monotonny~$\leftrightarrow$~inspirujący
      & j.w. \\
    \emph{Coherence} (spójność)
      & Relacja punktu do funkcji pokoju
      & Konfiguracja pokoju (aktywności, trudności) \\
    \emph{Extent} (rozległość)
      & Pośrednio: kontekst przestrzenny kadru
      & Analiza fotografii wnętrza przez VLM \\
    \bottomrule
  \end{tabular}
\end{table}

\begin{figure}[htbp]
  \centering
  \Todo{wstaw plik: figures/ch02/ch02-prs-mood-grid.png --- zrzut interfejsu mood grid z osiami X/Y oraz punktami pomiaru}
  \caption[Mapa nastroju PRS w interfejsie IDA: operacjonalizacja stanów ideal, current i target. Opracowanie własne na podstawie artefaktu; rama teoretyczna: Hartig i in.]{Mapa nastroju PRS w interfejsie IDA: operacjonalizacja stanów ideal, current i target.
    Opracowanie własne na podstawie artefaktu \cite{idaPlatform2026}; rama teoretyczna: \textcite{hartig1997prs}.}
  \label{fig:ch02:prs-mood-grid}
\end{figure}

\section{Projektowanie biofiliczne we wnętrzach}
\label{ch02:sec:biophilia}

Stephen Kellert, Judith Heerwagen i Martin Mador definiują projektowanie biofiliczne jako holistyczną strategię architektoniczną odpowiadającą na wrodzoną potrzebę więzi człowieka ze światem przyrody \footcite{kellert2008biophilic}. Zasada ta realizuje się nie tylko poprzez obecność żywych roślin, lecz przede wszystkim poprzez dostęp do światła dziennego, zastosowanie naturalnych, nieprzetworzonych materiałów oraz form biomorficznych.

W platformie IDA przełożyłem biofilię na czterostopniową skalę nasycenia wizualnego: od wnętrza całkowicie zneutralizowanego po przestrzeń o dominującym udziale zieleni, światła słonecznego i faktur organicznych.
Wybór ten bezpośrednio modyfikuje brief przekazywany modelowi generatywnemu, wpływając na gęstość roślinności, paletę materiałową oraz temperaturę barwną światła.

\begin{figure}[htbp]
  \centering
  \Todo{wstaw plik: figures/ch02/ch02-biophilia-dosage.png --- cztery poziomy nasycenia biofilicznego w interfejsie IDA}
  \caption[Stopnie nasycenia biofilicznego jako wybór wizualny w IDA. Opracowanie własne; rama teoretyczna: Kellert, Heerwagen i Mador.]{Stopnie nasycenia biofilicznego jako wybór wizualny w IDA.
    Opracowanie własne; rama teoretyczna: \textcite{kellert2008biophilic}.}
  \label{fig:ch02:biophilia}
\end{figure}

\section{Relacja człowieka z przestrzenią a konfiguracja pokoju}
\label{ch02:sec:peo}

Projektowanie wnętrza mieszkalnego jest procesem negocjacji między indywidualną tożsamością mieszkańca, jego codziennymi aktywnościami a zastanymi ograniczeniami budowlanymi.
W module konfiguracji pokoju (Room Setup) IDA zbiera kluczowe parametry tej relacji: przeznaczenie pomieszczenia, katalog dominujących aktywności, punkty tarcia (bariery użytkowe) oraz hierarchię wartości badaną metodą drabinowania (rozdz.~\ref{ch03:sec:laddering}).

Połączenie nastroju emocjonalnego (PRS) z programem funkcjonalnym chroni proces personalizacji przed degradacją do roli powierzchownej stylizacji. Dzięki temu generowany obraz odnosi się do realnego sposobu zamieszkiwania pokoju. \textcolor{orange}{W ujęciu poetyki przestrzeni Gastona Bachelarda dom i mieszkanie nie są bowiem obojętnym geometrycznie pudełkiem, lecz pierwszym wszechświatem człowieka i intymnym schronieniem wyobraźni, w którym proporcje wnętrza, światło i faktura materii przechowują pamięć zamieszkiwania \footcite{bachelard1975wyobraznia}.}

\section{Implikacje dla syntezy obrazu w IDA}
\label{ch02:sec:implikacje}

Z ram teoretycznych ART, SRT, skali PRS oraz biofilii wyłania się spójny zestaw wytycznych zasilających proces generowania koncepcji (tab.~\ref{tab:ch02:inputs-brief}).
Tabela ta ilustruje, w jaki sposób humanistyczne kategorie doświadczenia przestrzennego zostają przetłumaczone na konkretne sygnały wejściowe wieloźródłowej macierzy generacji w IDA.

\begin{table}[htbp]
  \centering
  \caption{Translacja wymiarów psychologiczno-przestrzennych na parametry wytycznych w macierzy IDA.}
  \label{tab:ch02:inputs-brief}
  \begin{tabular}{@{}>{\raggedright\arraybackslash}p{0.26\textwidth}>{\raggedright\arraybackslash}p{0.32\textwidth}>{\raggedright\arraybackslash}p{0.34\textwidth}@{}}
    \toprule
    Wymiar doświadczenia & Parametr w interfejsie IDA & Powiązane źródło macierzy \\
    \midrule
    Luka nastrojowa (PRS)
      & Odległość \texttt{current}$\rightarrow$\texttt{target}
      & Wariant \texttt{mixed\_functional} oraz \texttt{mixed} \\
    Nasycenie biofiliczne
      & Poziom 0--3 na skali wizualnej
      & Warianty \texttt{explicit}, \texttt{mixed}, \texttt{mixed\_functional} \\
    Funkcja i rytuały domowe
      & Aktywności, punkty tarcia, drabina wartości
      & Wariant \texttt{mixed\_functional} \\
    Preferencje estetyczne
      & Dyferencjał semantyczny, selekcja kart
      & Warianty \texttt{explicit} oraz \texttt{implicit} \\
    Profil osobowości
      & Kwestionariusz IPIP-NEO-120
      & Wariant \texttt{personality} \\
    Inspiracja wizualna
      & Zdjęcie referencyjne użytkownika
      & Wariant \texttt{inspiration\_reference} \\
    \bottomrule
  \end{tabular}
\end{table}

Psychologia środowiskowa w niniejszej pracy nie stanowi zatem celu samego w sobie, lecz staje się dyscyplinującą ramą projektową. Prowadzi od badania potrzeb regeneracyjnych i funkcji pokoju do kontrolowanej syntezy obrazu wnętrza, której techniczne mechanizmy i próby w ComfyUI omawia rozdział~\ref{ch04:top}.
